\subsection*{Retaining the archimedean energy in a one-sided estimate}

Choose a literal Fourier tail on which $a_n:=-A_n>0$ and write
$D=\operatorname{diag}(a_n)$ there. By
\eqref{eq:v114-weil-decomposition}, the compressed closed form has
$\mathsf W=D+R$, with $R$ bounded self-adjoint. It contains the
actual signed prime and pole terms and the archimedean off-diagonal;
no arithmetic diagonal is absorbed into $D$. Put
\begin{equation}\label{eq:v120-weighted-remainder}
 \mathcal B=D^{-1/2}R D^{-1/2}.
\end{equation}
The operator $\mathcal B$ is compact self-adjoint, since
$a_n\sim\log|n|$ at the fixed window. The form identity is
\[
 q_{\mathsf W}[f]
 =\ip{(I+\mathcal B)D^{1/2}f}{D^{1/2}f}.
\]
Thus the exact positivity target is $\mathcal B\succeq-I$;
the stronger two-sided condition $\|\mathcal B\|\le1$ can fail
because of a harmless positive direction.

\begin{proposition}[The scale of an approximate weighted sign]
\label{prop:v120-weighted-error}
Suppose $D\succeq d_0I>0$, $R\succeq-CI$, $C\ge0$, and
$D^{-1/2}RD^{-1/2}\succeq-(1+\eta)I$, $\eta\ge0$.
Then, on the entire form domain,
\begin{equation}\label{eq:v120-weighted-error}
 D+R\succeq-\frac{\eta C}{1+\eta}I.
\end{equation}
The constant is sharp. In particular, a dimensionless error
$\eta_\lambda\to0$ is not by itself an ordinary lower error
$o(1)$ when $C_\lambda$ grows.
\end{proposition}
\begin{proof}
The hypotheses give $D+R\succeq-\eta D$ and
$D+R\succeq D-CI$. Add these inequalities with weights
$1/(1+\eta)$ and $\eta/(1+\eta)$ respectively. For $C>0$,
the scalar example $D=C/(1+\eta)$, $R=-C$ attains equality.
Taking $C_\lambda=\lambda$ and $\eta_\lambda=1/\lambda$
shows why the missing scale factor cannot be omitted. A sufficient
conversion requires $\eta_\lambda C_\lambda\to0$, or a more
precise directional energy estimate.
\end{proof}

There is also a limitation on how the compact weighted remainder can
be estimated. A Hilbert--Schmidt budget for its whole infinite matrix
is unavailable, despite the valid finite-column residual Grams above.

\begin{proposition}[Weighted compactness without a square-summable matrix]
\label{prop:v120-not-schatten}
At any fixed $\lambda>\sqrt2$, let $D$ be a positive diagonal on a
fixed Fourier tail with $a_n\sim\log|n|$. The compact operator
$D^{-1/2}\mathsf T_{\rm pr}D^{-1/2}$ belongs to no Schatten class
$\mathcal S_p$, $0<p<\infty$. For the exact archimedean diagonal
$a_n=-A_n$, the same assertion holds for
$\mathcal B$ in \eqref{eq:v120-weighted-remainder}.
\end{proposition}
\begin{proof}
The prime diagonal is the nonzero finite trigonometric sequence
\[
 \tau_n=2\sum_{1<m<\lambda^2}
 \frac{\Lambda(m)}{\sqrt m}
 \left(1-\frac{\log m}{L}\right)
 \cos\frac{2\pi n\log m}{L}.
\]
Collect equal unit-circle frequencies to write
$\tau_n=\sum_{z\in F}c_z z^n$, with all $c_z>0$.
Finite geometric sums give
\[
 \frac1X\sum_{n=1}^X|\tau_n|^2\longrightarrow
 \sum_{z\in F}c_z^2>0.
\]
Boundedness then implies that $|\tau_n|$ exceeds a positive
constant on a set of positive lower density. Consequently, for each
$p\ge1$, $\sum_n|\tau_n/a_n|^p=\infty$: the partial sum on
that set is bounded below by a constant times $X/(\log X)^p$.
For self-adjoint $K\in\mathcal S_p$, Jensen's inequality in its
spectral measure gives
$\sum_n|\ip{KU_n}{U_n}|^p\le\operatorname{Tr}|K|^p$.
This is a contradiction. The inclusions
$\mathcal S_p\subset\mathcal S_1$ for $0<p<1$ settle the
remaining exponents. For the complete remainder its diagonal is
$(-\tau_n+O_\lambda(n^{-2}))/a_n$, since the exact pole diagonal
has this decay and the archimedean remainder has zero diagonal.
The same argument applies. None of this invalidates the separated
finite-column moment bounds, whose off-diagonal factor
$(n-m)^{-1}$ is square summable in the remote row index.
\end{proof}

\begin{proposition}[Sharp fixed-window decay of the weighted prime tail]
\label{prop:v120-weighted-tail-asymptotic}
Under the diagonal assumptions above, let $Q_J$ remove $|n|\le J$.
At each fixed window,
\begin{equation}\label{eq:v120-weighted-tail-asymptotic}
 \lim_{J\to\infty}(\log J)
 \|Q_JD^{-1/2}\mathsf T_{\rm pr}D^{-1/2}Q_J\|
       =\|\mathsf T_{\rm pr}\|.
\end{equation}
Subtracting any fixed compact operator from $\mathsf T_{\rm pr}$
before weighting leaves this limit unchanged. No uniformity in
$\lambda$ is asserted.
\end{proposition}
\begin{proof}
The upper bound is
$\|\mathsf T_{\rm pr}\|/\inf_{|n|>J}a_n$.
For the lower bound, fix $\epsilon>0$ and choose a unit Fourier
polynomial $f$, supported on $|m|\le M$, with
$|\ip{\mathsf T_{\rm pr}f}{f}|>\|\mathsf T_{\rm pr}\|-\epsilon$.
Returns of the finite vector of shift phases to any fixed neighborhood
of the identity are relatively dense. Indeed its cyclic closure is a
compact group; finitely many integer translates of that neighborhood
cover it. The largest difference between the corresponding translation
indices bounds the gap between return times. This also covers rational
phase relations.

Choose the neighborhood so that the conjugation error in
\eqref{eq:v119-prime-recurrence} is below $\epsilon$.
For every large $J$ there is therefore a return
$k\in[J+M+1,J+M+1+K]$, with $K$ independent of $J$.
Set $g=M_kf$ and $x=D^{1/2}g$. Both are supported outside $J$,
$\|x\|^2=(1+o(1))\log J$, and the weighted Rayleigh numerator
at $x$ equals $\ip{\mathsf T_{\rm pr}g}{g}$, whose absolute value
is greater than $\|\mathsf T_{\rm pr}\|-2\epsilon$.
Let $J\to\infty$ and then $\epsilon\downarrow0$.
For compact $C$ the change in the left side is at most
\[
 \frac{\log J}{\inf_{|n|>J}a_n}\|Q_JCQ_J\|\longrightarrow0.
\]
The return-length bound $K$ depends on the window and the tolerance.
Thus one may not combine \eqref{eq:v120-weighted-tail-asymptotic}
with $\|\mathsf T_{\rm pr}\|\sim\lambda$ to infer a bound on
an arbitrary joint scale $J=J(\lambda)$.
\end{proof}

\begin{proposition}[A signed energy bound at the larger window $\lambda=4$]
\label{prop:v120-lambda4-tail}
Let $L=2\log4$, retain the exact diagonal $a_n=-A_n$, and let
$Q_{16}$ be the literal projection onto $|n|>16$. For every complex
vector in its complete closed form domain,
\begin{equation}\label{eq:v120-lambda4-tail}
 QW_4(f,f)\ge10^{-8}\sum_{|n|>16}a_n|f_n|^2,
 \qquad f=Q_{16}f.
\end{equation}
Every weight in this sum is positive. This is an infinite-tail
computer-assisted bound, not a claim about the remaining low modes
or about an unbounded family of windows.
\end{proposition}
\begin{proof}
Apply the directional verified-solve proof to
$\mathsf W_4-cD$, $c=10^{-8}$, on $Q_{16}$.
Only the archimedean diagonal is changed, to $(1-c)a_n$;
the exact off-diagonal coefficients $b_n$, both pole signs, and every
prime power $2,3,4,5,7,8,9,11,13$ remain. The full-length $m=16$
translation is zero. The positive physical weight
$\phi(x)=e^{-x/2}+e^{-(L-x)/2}$ gives the certified Schur bound
\begin{equation}\label{eq:v120-lambda4-prime-bound}
 \|\mathsf T_{\rm pr}\|
 \le M_\phi<4.624042316766478.
\end{equation}
To verify it, set $u=e^x$. On each interval between
$1,16,m,16/m$ the ratio $(\mathsf T_{\rm pr}\phi)/\phi$
is a linear-fractional function of $u$, hence monotone or constant.
Every one-sided endpoint value is evaluated with outward rounding;
the chosen upper endpoint is checked against every candidate.

Use $E_L,C_L,h_L$ from \eqref{eq:v115-arch-errors}, with
$C_L=1$, $h_L=9/16$ and $R_L=64/255$. For $n>N$ the proved
remote inverse bound applies with
\[
 g_n=(1-c)\{\log(n/L)-E_L(2\pi n/L)\}-\kappa,
\]
where $\kappa=M_\phi+2C_L/t_*$ in the even sector and
$\kappa=M_\phi+2C_L/t_*+\pi/2+4h_LL/(\pi^2N)$ in the odd
sector, $t_*=2\pi(N+1)/L$. These weights are positive and
increasing for the choices below. Also the same diagonal estimate
at $n=17$ proves positivity of $a_n$ on the whole claimed tail.

The finite head now consists of $17\le n\le N$ in the normalized
parity basis; it does not include the omitted physical low modes.
The solve witness has support $N<n\le M$. With
$G=I-Z$, form the exact signed $K_Z$ and every residual row up to
$J$, and use the moment majorant $U_J$ beyond $J$. The lower matrix is
\begin{equation}\label{eq:v120-lambda4-lower}
 \mathcal L=K_Z-
  \sum_{N<n\le J}\frac{R_n^*R_n}{g_n}-\frac{U_J}{g_{J+1}}.
\end{equation}
The same full-index parity conversion and both infinite tails as in
Proposition~\ref{prop:v116-remote-moment} are used. Bounding its
geometric-remainder constant by including the absent indices
$1,\ldots,16$ is conservative. The certificate parameters are
\begin{center}
\begin{tabular}{lrrrrr}
\toprule
sector&head size&$N$&$M$&$J$&moment order\\
\midrule
even&496&512&1024&4096&48\\
odd&1520&1536&2048&4096&64\\
\bottomrule
\end{tabular}
\end{center}
Both $Z$ and an inverse-Cholesky congruence witness $C$ are frozen
to exact dyadic entries. Their numerical method of construction has
no proof role: outward interval evaluation verifies every row of
$C\mathcal L C^*$ has positive diagonal minus the sum of the
absolute off-diagonal entries. Thus $C\mathcal L C^*\succ0$
and $\mathcal L\succ0$. Square completion with the positive
remote operator proves \eqref{eq:v120-lambda4-tail} separately
in each parity sector and hence for arbitrary complex vectors.

The saved witnesses, analytic series radii, all finite residual rows
and the infinite moment bound are reproduced by
\texttt{assembly\_general.py} and
\texttt{certify\_weighted\_tail.py}. Replaying the same witnesses
at higher precision verifies the same strict inequalities.
These are internal computer-assisted certificates, not independent
external verification. No zero locations or RH assumptions enter them.
\end{proof}

\begin{proposition}[The two-sided weighted contraction fails at this cutoff]
\label{prop:v120-norm-counterexample}
For the same exact $D$ and $Q_{16}$ at $\lambda=4$,
\[
 \|D^{-1/2}Q_{16}(\mathsf W_4-D)Q_{16}D^{-1/2}\|>1.
\]
This does not contradict Proposition~\ref{prop:v120-lambda4-tail}.
\end{proposition}
\begin{proof}
The archived exact dyadic even vector $v$, supported on
$17\le|n|\le256$, has the rigorously enclosed quotient
\[
 2.05254409540345<\frac{QW_4(v,v)}{\ip{Dv}{v}}
                         <2.05254409540346.
\]
Its positive numerator excess over $2\ip{Dv}{v}$ exceeds
$0.0525$ in the saved normalization. The same analytic matrix
assembly, independent of the floating eigenvector that suggested
$v$, proves these interval inequalities. Testing the weighted
remainder on $D^{1/2}v$ gives a Rayleigh value greater than one.
This is a positive direction, not a negative Weil direction.
\end{proof}

The larger-window tail sign is now established, but full positivity
at $\lambda=4$ still requires the signed effective matrix on
$E_{16}$:
\[
 F-B^*T^{-1}B,\qquad
 F=P_{16}\mathsf W_4P_{16},\quad B=Q_{16}\mathsf W_4P_{16},
 \quad T=Q_{16}\mathsf W_4Q_{16}.
\]
It has $17$ even and $16$ odd coordinates. It is not the raw finite
Weil matrix: the inverse-tail correction must remain.
The next growing-window estimate must also satisfy the ordinary-error
scale in Proposition~\ref{prop:v120-weighted-error}.
A related recent preprint~\cite{Zhu2026Windows} studies certified
compact-window positivity and a pointwise prime-comb envelope barrier.
Our conclusion here concerns the actual Fourier-tail form and does
not use that preprint's computational certificates. Neither fixed-window
result supplies a uniform signed arithmetic estimate.
